\section{Discussion}\label{sec:discussion-new}

The main tools used here are classical: companion matrices, rational generating functions, finite Hankel rank, Weierstrass reduction, Riemann--Hurwitz, local Puiseux expansions, and the Beraha--Kahane--Weiss description of limiting zeros.
What is specific to this paper is that these tools can be pushed to completion for the concrete quartic family \(\Pi(\lambda,y)\).
For this family we obtain the explicit factorization of the ramification polynomial, the exact real-branch classification, the real-spectral phase diagram, the global dominant-sheet description on \(y>0\), the local endpoint zero quantization at \(y=0\), and the exact degree law.

The scope of the present analysis is limited but precise.
All results concern the quartic companion family defined in Section~\ref{sec:model}.
No microscopic derivation of this family is claimed here, nor is one needed for the spectral statements proved above.
Accordingly, terms such as dominant sheet, branch point, and \(f_{\mathrm{dom}}(y)=\log r(y)\) are used purely in the spectral sense internal to the quartic family itself.

From this viewpoint, the main structural point is the separation between the full ramification locus and the dominant-modulus envelope.
At \(y=0\), the branch collision occurs on the dominant sheet over \(y>0\) and therefore controls the first edge zeros.
At \(y=1\) and \(y=y_{\mathrm{sub}}\), by contrast, the collisions are subdominant and do not create the same boundary behavior.
The exact real-spectral phase diagram makes this distinction visible globally, while the local normal form of Section~\ref{sec:physical-sheet} explains the resulting odd-integer edge quantization.

At the opposite end of the parameter range, the large-\(y\) expansion is governed by a double tangency at projective infinity.
This gives a second exact mechanism, distinct from the finite branch-point analysis, for the quarter-power splitting of the dominant root.
Together with the degree law of Section~\ref{sec:support-law}, it shows that the top-degree sector is rigid in a precise algebraic sense.

Two extensions appear natural.
One is to derive the quartic family from an explicit microscopic model and thereby recover a fully internal statistical interpretation of the coefficients.
The other is to push the balanced-family viewpoint beyond the normalization step and determine how much of the present branch classification and dominant-envelope analysis persists across that larger parameter space.
